% ============================================================
% CHAPTER 9 — Project 6: Web Frontend with React and Authentication
% ============================================================
\chapter{Project 6: Web Frontend with React and Authentication}\index{React}\index{Tailwind CSS}\index{Vite}\index{SKILL.md}

\section*{What You Will Build}
A complete React frontend application that:
\begin{itemize}
  \item Connects to the Notes API backend
  \item Allows login with Google and GitHub
  \item Shows a dashboard with the user's notes
  \item Supports full CRUD operations
  \item Has protected routes
  \item Uses a consistent design system with Tailwind~CSS
\end{itemize}

\textbf{Estimated time}: 60--90 minutes\\
\textbf{Prerequisite}: Notes API backend with authentication (Chapters~6--8)

\begin{notebox}
\textbf{Tooling Box --- Stack chosen for this example.}
\begin{itemize}
  \item \textbf{Framework:} React~19 with Vite
  \item \textbf{Styling:} Tailwind~CSS~4
  \item \textbf{Routing:} React Router~7
  \item \textbf{HTTP Client:} Axios
\end{itemize}
\textbf{Equivalent alternatives:} Vue.js/Nuxt, Svelte/SvelteKit, Angular, Next.js, Astro. The \textbf{pattern} (components, protected routing, state management, API calls) is the same in any modern frontend framework. The \texttt{SKILL.md} file you will create in this chapter works with any stack --- you only need to adapt the conventions.
\end{notebox}

% ---------------------------------------------------------
\section{Frontend Project Setup}

\begin{praticabox}[title={Hands-On --- Create the React project}]
Open a terminal in the parent directory of the backend:
\begin{lstlisting}[language=bash]
cd ..
npm create vite@latest notes-frontend -- --template react
cd notes-frontend
npm install
npm install react-router-dom axios
npm install -D tailwindcss @tailwindcss/vite
\end{lstlisting}
Now you have two projects side by side:
\begin{lstlisting}[language={}]
workspace/
+-- notes-api/          <- Backend (Chapters 6-8)
+-- notes-frontend/     <- Frontend (this chapter)
\end{lstlisting}
\end{praticabox}

% ---------------------------------------------------------
\section{Agent Skill for the Frontend}

This is the moment to introduce your first \textbf{SKILL.md} --- a document that gives the AI specialized expertise.

\begin{praticabox}[title={Hands-On --- Create the file \texttt{SKILL.md}}]
\end{praticabox}

\begin{lstlisting}[language={},title={\texttt{SKILL.md} --- React Frontend Developer}]
# React Frontend Developer Skill

## Design Principles
- Mobile-first
- Small components: each component does ONE thing
- No business logic in UI components
- Accessibility: interactive elements with keyboard support

## Component Structure
src/
  components/
    ui/         <- Button, Input, Card, Modal
    layout/     <- Header, Footer, PageContainer
    notes/      <- NoteCard, NoteForm, NoteList
    auth/       <- LoginButton, ProtectedRoute, UserMenu
  pages/
    LoginPage.jsx
    DashboardPage.jsx
    NoteDetailPage.jsx
    NotFoundPage.jsx
  hooks/
    useAuth.js     <- Hook for authentication
    useNotes.js    <- Hook for notes CRUD
    useApi.js      <- Hook for API calls
  context/
    AuthContext.jsx
  services/
    api.js         <- Configured Axios instance

## React Conventions
- Components: PascalCase, one file per component
- Hooks: camelCase with the "use" prefix
- Event handlers: prefix "handle"
- Only functional components with hooks
- NEVER make fetch/axios calls in components

## Constraints
- DO NOT use localStorage for JWT tokens
- DO NOT hardcode the backend URL. Use VITE_API_URL
- Every page must have a loading state and an error state
- Forms with client-side validation before submit
\end{lstlisting}

\begin{notebox}
The file \texttt{SKILL.md} works like \texttt{\_CONTEXT.md} but is \textbf{specialized}. It is loaded by the AI only when the type of work requires it (Progressive Disclosure). This avoids overloading the context.
\end{notebox}

% ---------------------------------------------------------
\section{The Frontend Context}

\begin{praticabox}[title={Hands-On --- Create \texttt{\_CONTEXT.md} in the frontend}]
\end{praticabox}

\begin{lstlisting}[language={},title={\texttt{\_CONTEXT.md} for Notes Frontend}]
# Project: Notes Frontend

## Purpose
React frontend for the Notes application. It connects to the
Notes API backend (http://localhost:3000).

## Technologies
React 18 + Vite, React Router 6, Tailwind CSS 4, Axios,
React Context for global state (auth).

## Authentication Flow
1. User clicks "Login with Google/GitHub"
2. Redirect to the backend /api/auth/google
3. Backend handles OAuth, returns JWT as a cookie
4. Frontend calls /api/auth/me for the profile
5. AuthContext saves the user in memory
6. Subsequent API calls include the cookie

## Pages
/              LoginPage (if not authenticated)
/dashboard     DashboardPage (list notes + creation)
/notes/:id     NoteDetailPage (detail + edit)
*              NotFoundPage

## Constraints
- DO NOT use localStorage for JWT tokens
- Axios with withCredentials: true
- Vite proxy to avoid CORS issues
\end{lstlisting}

% ---------------------------------------------------------
\section{Generating the Frontend}

\begin{praticabox}[title={Hands-On --- Generate the application}]
In Copilot Agent Mode:
\begin{lstlisting}[language={}]
Read _CONTEXT.md and SKILL.md. Implement the entire
React frontend.

Order:
 1. Configure Tailwind CSS with Vite
 2. Create src/services/api.js (configured Axios)
 3. Create src/context/AuthContext.jsx
 4. Create basic UI components (Button, Input, Card)
 5. Create layout (Header with UserMenu, PageContainer)
 6. Create LoginPage with Google and GitHub buttons
 7. Create ProtectedRoute
 8. Create src/hooks/useNotes.js
 9. Create DashboardPage
10. Create notes components (NoteCard, NoteForm)
11. Create NoteDetailPage
12. Configure React Router in App.jsx
13. Configure the Vite proxy
14. Create .env with VITE_API_URL
\end{lstlisting}
\end{praticabox}

% ---------------------------------------------------------
\section{Vite Proxy Configuration}

To avoid CORS issues during development:

\begin{lstlisting}[language=JavaScript,title={\texttt{vite.config.js}}]
import { defineConfig } from 'vite';
import react from '@vitejs/plugin-react';
import tailwindcss from '@tailwindcss/vite';

export default defineConfig({
  plugins: [react(), tailwindcss()],
  server: {
    port: 5173,
    proxy: {
      '/api': {
        target: 'http://localhost:3000',
        changeOrigin: true,
        secure: false
      }
    }
  }
});
\end{lstlisting}

\begin{tipbox}
With this proxy, in the frontend you can call \texttt{/api/notes} instead of \texttt{http://localhost:3000/api/notes}. Vite automatically forwards the requests to the backend.
\end{tipbox}

% ---------------------------------------------------------
\section{Testing the Application}

\begin{praticabox}[title={Hands-On --- Start-up}]
Open \textbf{two} terminals:

\textbf{Terminal 1 --- Backend:}
\begin{lstlisting}[language=bash]
cd notes-api && npm run dev
\end{lstlisting}

\textbf{Terminal 2 --- Frontend:}
\begin{lstlisting}[language=bash]
cd notes-frontend && npm run dev
\end{lstlisting}

Open in the browser: \texttt{http://localhost:5173}
\end{praticabox}

\begin{center}
\begin{tabularx}{\textwidth}{l X}
\toprule
\textbf{Feature} & \textbf{Test} \\
\midrule
Login Page & You see the Google and GitHub buttons \\
OAuth Google & Click $\rightarrow$ redirect $\rightarrow$ login $\rightarrow$ dashboard \\
Dashboard & Note list (empty at first) \\
Create note & Fill in the form $\rightarrow$ appears in the list \\
Detail & Click a note $\rightarrow$ detail page \\
Edit & Edit title/content $\rightarrow$ save $\rightarrow$ verify \\
Delete & Click delete $\rightarrow$ confirm $\rightarrow$ disappears \\
Logout & Click logout $\rightarrow$ redirect to login \\
Route protection & Go to /dashboard without login $\rightarrow$ redirect \\
User Menu & User name and avatar in the header \\
\bottomrule
\end{tabularx}
\end{center}

\begin{checkpointbox}
If you can log in, create a note, edit it, and delete it, the frontend is complete.
\end{checkpointbox}

% ---------------------------------------------------------
\section{Design Iterations}

\begin{praticabox}[title={Hands-On --- Improve the appearance}]
\begin{lstlisting}[language={}]
Improve the design:
1. EmptyState when there are no notes
2. Hover animations on cards (scale-[1.02])
3. Tags as colored badges on the NoteCards
4. Search field in the dashboard
5. Relative creation date ("2 hours ago", "yesterday")
\end{lstlisting}
\end{praticabox}

\begin{praticabox}[title={Hands-On --- Responsive design}]
\begin{lstlisting}[language={}]
Verify that the interface is responsive:
1. Notes grid: 1 col mobile, 2 tablet, 3 desktop
2. Hamburger header on mobile
3. Full-width form on mobile
4. Buttons with touch target >= 44px
\end{lstlisting}
\end{praticabox}

% ---------------------------------------------------------
\section*{Summary}

\begin{center}
\begin{tabularx}{\textwidth}{l X}
\toprule
\textbf{Aspect} & \textbf{Detail} \\
\midrule
Framework & React~18 + Vite \\
Styling & Tailwind~CSS~4 \\
Auth & OAuth~2.0 via backend proxy \\
State & React Context (AuthContext) \\
Pages & Login, Dashboard, NoteDetail, 404 \\
First SKILL.md & react-frontend-developer \\
Time & $\sim$60--90 minutes \\
\bottomrule
\end{tabularx}
\end{center}
